\chapter*{Streszczenie}
\addcontentsline{toc}{chapter}{Streszczenie}

W pracy zbadano możliwość użycia autokoderów do rozstrzygania koreferencji w polskich tekstach prawniczych. Przeanalizowano literaturę dotyczącą koreferencji neuronowej, polskich zasobów, przetwarzania języka naturalnego (natural language processing, NLP) w prawie, autokoderów oraz dużych modeli językowych (large language models, LLM). Zaprojektowano wspólny potok danych i ewaluacji oraz dwa warianty: odszumiający autokoder (denoising autoencoder, DAE) reprezentacji wzmianek z głowicą parową i U-Net, czyli konwolucyjny enkoder--dekoder z połączeniami skrótowymi, segmentujący wielokanałową macierz relacji. Dodano kontrolę neuronową bez rekonstrukcji, baseline’y klasyczne oraz selektor hybrydowy kierujący niepewne przypadki do LLM. Zaimplementowano trening sterowany formatem YAML (YAML Ain't Markup Language), deterministyczne seedy, checkpointy, oficjalny scorer projektu Coreference in Universal Dependencies (CorefUD) i pomiar kosztu. W zredukowanej próbie syntetycznej kontrola uzyskała średnio 0,94632 w metryce Conference on Natural Language Learning F1 (CoNLL F1), DAE 0,85544, a U-Net 0,89270. Wyniki służą wyłącznie weryfikacji integracji i nie rozstrzygają jakości na polszczyźnie ani w domenie prawnej. Nie wykonano eksperymentów LLM i testu na ręcznie anotowanych orzeczeniach z powodu braku danych, klucza do interfejsu programowania aplikacji (application programming interface, API) i końcowego audytu prywatności. Tezy o przewadze domenowego autokodera nie potwierdzono ani nie obalono w zaplanowanym sensie. Rezultatem jest działająca infrastruktura oraz jawny, falsyfikowalny protokół dalszej ewaluacji.

\noindent\textbf{Słowa kluczowe:} koreferencja, autokoder, teksty prawnicze, język polski, U-Net, duże modele językowe, CorefUD.

\chapter*{Abstract}
\addcontentsline{toc}{chapter}{Abstract}

This thesis investigates whether autoencoders can support coreference resolution in Polish legal texts. Literature on neural coreference, Polish resources, legal natural language processing (NLP), autoencoders, and large language models (LLMs) was reviewed. A shared data and evaluation pipeline was designed together with two variants: a denoising autoencoder (DAE) for mention representations followed by a pairwise scorer, and a U-Net convolutional encoder--decoder segmenting a multi-channel relation matrix. A neural control without reconstruction, classical baselines, and a hybrid selector sending uncertain cases to an LLM were added. The implementation provides YAML (YAML Ain't Markup Language) configuration files, deterministic seeds, checkpoints, the official Coreference in Universal Dependencies (CorefUD) scorer, and cost measurements. In a reduced synthetic integration run, the neural control achieved a mean Conference on Natural Language Learning F1 (CoNLL F1) of 0.94632, compared with 0.85544 for the DAE and 0.89270 for the U-Net. These values validate the technical pipeline only; they do not measure performance on Polish or legal language. The LLM experiments and evaluation on manually annotated judgments were not run because the required data, application programming interface (API) credentials, and final privacy audit were unavailable. Consequently, the claim that domain-trained autoencoders improve legal coreference was neither confirmed nor rejected under the predefined protocol. The main outcomes are a working experimental infrastructure, explicit limitations, and a falsifiable procedure for completing the comparison on frozen natural-language datasets.

\noindent\textbf{Keywords:} coreference resolution, autoencoder, legal texts, Polish language, U-Net, large language models, CorefUD.
